% English translation of capitulos/04c-qp3.tex, made by Claude Code (Claude Fable 5.2) on 24 November 2026.
% Pending review by the author. When the Portuguese file changes, update this one (scripts/verificar_traducao.py).
\subsection{RQ3: taxonomy and decision-support table}\label{sec:qp3}

This subsection brings together the practical contribution of the work: a taxonomy of the techniques
extracted and a decision-support table, in the version revised after the examination in the scenarios. Both are based on 332 records
and 805 comparisons extracted by AI and checked by the author in a sample (Section~\ref{sec:conferencia}). The counts come from
\texttt{scripts/sintetizar.py}, and the rows of the table, from the partial syntheses.

\subsubsection{Taxonomy}

The taxonomy is multi-label: a study may have more than one task, and a technique, more
than one mechanism. The records register five dimensions:
\begin{itemize}
  \item \textit{Task}: generation (247 studies), debugging (62), translation (27) and
    refactoring (20); 22 studies have more than one task.
  \item \textit{Mechanism}: eleven labels (instruction, context, constraints, examples,
    decomposition, reasoning, feedback, sampling, multi-agent, prompt optimization and
    other). The record separates the complete package, with all the mechanisms present in the prompt
    of the technique, from the mechanisms that the technique adds to the base prompt;
    Table~\ref{tab:mecanismos} counts only those added.
  \item \textit{Role}: focal, when the technique is the object of the study, or comparator, when
    it serves as a base or as an alternative to another technique.
  \item \textit{Type of counterfactual}: package (whole technique against the condition without
    it), ingredient (an isolated component), paired contrast, cumulative (successive
    rounds, without a parallel condition) or none. Among the focal techniques, there are 373 of
    package, 212 of ingredient, 3 of paired contrast, 13 cumulative and 111 without a
    counterfactual (unit: technique). Counting also the components of the techniques, 191 of the
    332 studies have at least one isolated ingredient or paired contrast.
  \item \textit{Conditions}: factors on which the results depend, according to the syntheses.
\end{itemize}

The syntheses point to the following recurring conditions:
\begin{itemize}
  \item available executor, tests or compiler, on which feedback and
    candidate selection depend;
  \item number of calls and of candidates, rarely equalized between the conditions;
  \item size and family of the model, with sign reversals within the same study;
  \item level of the task: function, class or repository;
  \item oracle information, such as perfect localization of the defect or APIs and solutions
    taken from the reference;
  \item components outside the scope of prompting (fine-tuning, trained retrievers,
    model selection), recorded in 303 of 332 studies;
  \item metric (textual similarity, tests or scanners) and domain, such as hardware and
    domain languages, in general with their own sets.
\end{itemize}

\subsubsection{Decision-support table}

Version 1 of the table consolidated the rows of the partial syntheses into 22, one per grouping
and task, and was frozen for the examination in the scenarios (Section~\ref{sec:exame}). After the
examination, the author decided to revise it. Version 2 (Tables~\ref{tab:apoio} and~\ref{tab:apoiob})
uses the same syntheses and has 31 rows, with five changes:
\begin{itemize}
  \item feedback was separated by type of verification, tests or compiler and
  analyser, in generation (G1 and G2) and in translation (T1 and T2);
  \item each row states the measured result, such as compilation, passing the tests,
  correctness of the repair or complexity;
  \item the condition makes explicit when the evidence assumes perfect localization or a test that
  reproduces the defect (D1, D3 and D6);
  \item the refactoring row was divided into specified type, API knowledge and
  examples (R1 to R3);
  \item each indication is marked as favourable, contrary, mixed, indirect or
  insufficient. Five new rows declare insufficient evidence: the distribution of a
  fixed budget of calls (G13), domain languages and small models (G14), repair
  without reproduction (D7), legacy language (T5) and repository scale (R5). The separations
  created G2, with indirect evidence, T2, insufficient, and R1 to R3.
\end{itemize}
In the mixed, contrary or insufficient rows, the guidance is not to recommend the technique
as a default. The strength is descriptive, without a score, and repeats the count of v1 when a row was
divided. v1 and v2 are in the repository (\path{evidencias/sinteses/}). The findings
below support the central rows.

In execution feedback in generation, the budget matters.
% AFIRMACAO: A91
With the number of attempts equalized, the multi-turn variants with feedback stayed at
2.4\% and 2.8\% of pass 1@3, against 2.9\% without iteration (Llama 3.0 8B), and at 6.9\% against
7.7\% (Llama 3.1 8B); SEIDR solved on average 11.5 problems of PSB2, against 12.2 for
sampling without repair (S479-C2; S479-C3; S479-C4; S645-C1).
% AFIRMACAO: A92
Without this equalization, five comparisons that isolate the loop and have absolute
values recorded a more favourable outcome with it, in the sense of each metric, and removing the iterative verification from CG-CG
reduced Pass@1 from 80.3\% to 70.1\% (S878-C1; S1047-C3; S1053-C3; S1077-C1;
S1086-C6; S1977-C2).
% AFIRMACAO: A93
In multi-agent systems, the gain comes with more calls: MapCoder and
LLMDebugger obtained 64.02\% and 67.07\% pass@1, against 39.63\% for \textit{zero-shot}, with 8.1
and 3.4 calls on average, against 1.0 (S1747-C3; S1747-C4).

% AFIRMACAO: A94
The reasoning instruction added to the generation prompt had a mixed result: in
12 comparisons from 12 studies, under heterogeneous conditions and metrics, 5 record an improvement and 7, a deterioration, a
descriptive count that is not equivalent to 12 isolated effects of the instruction (S124-C2; S207-C1; S309-C2; S319-C1; S491-C3; S508-C1;
S560-C2; S569-C1; S1024-C4; S1121-C1; S1229-C2; S1606-C4).
% AFIRMACAO: A95
Self-correction without execution, with Qwen2.5-1.5B, stayed below Simple-prompting: 38.2\%
(Self-Refine) and 38.8\% (SETS), against 44.6\% Accuracy@t2 (S759-C1; S759-C3).
% AFIRMACAO: A96
In the security instructions, 13 variants without iteration raised the vulnerable samples
of GPT-3.5-turbo from 5.74\% to values between 5.99\% and 7.82\% (S604-C1 to S604-C10;
S604-C13 to S604-C15), whereas, with GPT-3, the instructional variants reduced the
weaknesses counted by Bandit from 78 to values between 47 and 72 (S397-C2 to S397-C6).
Since model, metric and set differ, the two studies do not form a direct contrast.
% AFIRMACAO: A97
Oracle information produced the largest gains among the contrasts cited, without representing real use: the file
context added to the APIs taken from the reference raised the Pass@1 of GPT-4o Mini on RepoExec from 16.62\% to 37.75\%
(S670-C5), and the official solution raised the pass@1 of DeepSeek-Coder-V2 on MuPA from 11.84 to 38.79, the highest value among the conditions cited, whereas the self-generated hints stayed between 11.92 and 13.21 (S2145-C1; S2145-C3;
S2145-C5; S2145-C6).

% AFIRMACAO: A98
In repair, explicit reasoning worsened the result in 5 of 6 comparisons from 5
studies; the exception, the \textit{few-shot} CoT of ThinkRepair, also changes the selection of examples
(S275-C2; S354-C2; S1429-C1; S1429-C2; S1598-C2; S412-C2).
% AFIRMACAO: A99
Structural or bulky diagnostic information also made things worse: with Qwen2.5-Coder:7B on
QuixBugs, PDG and test file reduced the repair rate from 62.5\% to 55\% and
57.5\%, and, on Defects4J v1.2, prompts with stack trace localized 32 and 41 bugs, against 47
without it (S1674-C3; S1674-C4; S1678-C1; S1678-C2).
% AFIRMACAO: A100
Three repair comparisons with fixed examples recorded no improvement: the
overall\_score fell from 0.783 to 0.643, the plausible patches fell from 205 to 200, and
the failures fixed stayed at 11 against 11 (S1598-C1; S2057-C1; S1565-C4).

% AFIRMACAO: A101
In translation, explaining before translating reduced pass@1 from Python to JavaScript with
GPT-3.5 from 85.5\% to 84.8\%, 84.3\% and 83.2\%, and the two-step approaches reduced
the success rate on CodeNet from 28.6\% to values between 2.6\% and 11.2\% (S193-C1;
S193-C2; S193-C3; S836-C2; S836-C3; S836-C5; S836-C6).
% AFIRMACAO: A102
In refactoring, specifying the type of refactoring raised the recall of GPT-4 from 15.6\% to
52.2\%, and the TerraMod package raised the Pass@1 of Mistral-Large from 4\% to 76\%, both
from low baselines (S511-C2; S2122-C1).

% English translation of capitulos/04c2-tabela-apoio-v2.tex, made by Claude Code (Claude Fable 5.2) on 24 November 2026.
% The Portuguese file is generated by scripts/gerar_tabela_apoio.py from evidencias/sinteses/tabela-apoio-v2.csv;
% this file translates the generated .tex and is not regenerated by the script. Pending review by the author.
\providecommand*{\rr}{\raggedright\hspace{0pt}}
\begin{table*}[!tp]
\caption{Decision-support table, version 2: generation (G1--G12) and common rows (G13--G14)}\label{tab:apoio}
\centering
\fontsize{6.6}{7.4}\selectfont
\renewcommand{\arraystretch}{0.95}
\setlength{\tabcolsep}{2.5pt}
\renewcommand*{\rr}{\raggedright\hspace{0pt}}
\begin{tabular}{p{0.55cm}p{2.0cm}p{2.4cm}p{1.9cm}p{5.6cm}p{1.6cm}p{3.1cm}}
\hline
\rr\textbf{ID} & \rr\textbf{Grouping} & \rr\textbf{Condition and verification} & \rr\textbf{Measured result} & \rr\textbf{Evidence and indication} & \rr\textbf{Strength} & \rr\textbf{Caveats}\tabularnewline
\hline
\rr G1 & \rr Test or execution feedback in a loop & \rr Executable tests or executor; single-call base & \rr Pass@1; passing the tests & \rr \textit{Favourable; contrary with equalized budget.} Higher value with the loop, in a package (S150-C1) and isolated (S150-C2; S1977-C2). In the conditions extracted with the number of attempts equalized, no gain over sampling (S479-C2; S645-C1) & \rr approx. 37 (pkg. approx. 21; ingr. approx. 15), count of v1 for all types of feedback & \rr Calls not equalized in most; other ingredients in the package; passing the tests does not guarantee the requirement\tabularnewline
\rr G2 & \rr Compiler or static analyser feedback & \rr Compiler, IDE or security analyser, without functional tests & \rr Compilation; secure generation accuracy; validity & \rr \textit{Indirect.} Few studies and weak designs: compilable from 8 to 17 in cumulative rounds, without a parallel condition (Solidity; S2055-C1); security feedback in a package with a 32B model (S1065-C1); static report tied with self-review (S1763-C3) & \rr 3 studies & \rr Compiling does not measure functional correctness; a scanner does not demonstrate security\tabularnewline
\rr G3 & \rr Self-critique without execution & \rr The model itself critiques and rewrites, without executing & \rr Weaknesses flagged by scanner; accuracy; completion rate & \rr \textit{Mixed.} Fewer weaknesses with RCI (S397-C1; S604-C11); deterioration with a 1.5B model in one study (S759-C1); small gain (S1204-C2). Do not recommend as a default & \rr 6 (ingr. 1) & \rr One deterioration with 1.5B does not establish the behaviour of all small models; statistical null in S397\tabularnewline
\rr G4 & \rr Multi-agent with roles & \rr Tester, reviewer or evaluator; more calls, in a number that varies per arrangement & \rr Pass@1; compilation (S492) & \rr \textit{Favourable to the package.} Gains against a direct base (S344-C1; S2131-C1) and in ablations (S448-C3); deterioration in compilation in S492-C1 & \rr approx. 29 (pkg. approx. 29; ablation approx. 11) & \rr Arrangement against single agent only in S530; MapCoder with 8.1 calls on average (S1747-C3); fine-tuning of agents\tabularnewline
\rr G5 & \rr Sampling and selection & \rr Several candidates; selection by generated tests or value model & \rr Pass@1; creativity (S399) & \rr \textit{Favourable to selection; contrary to simple sampling.} Improvement in selection ablations (S416-C2; S724-C2; S1034-C2); deterioration with simple sampling (S399-C3) & \rr approx. 17 (ingr. approx. 9) & \rr Selection by the evaluation oracle is an upper bound (S399); trained components (S1034, S1592)\tabularnewline
\rr G6 & \rr Reasoning instruction (CoT) & \rr The same model generates the reasoning & \rr Pass@1 & \rr \textit{Mixed.} 5 improvements and 7 deteriorations in 12 comparisons (S319-C1; S1229-C2; S569-C1; S1121-C1). Do not recommend as a default & \rr 12 (ingr. 11; pkg. 1) & \rr Non-significant null (S560); token ceiling (S1606); reversals between models\tabularnewline
\rr G7 & \rr Planning or decomposition before the code & \rr Two or more steps; function level & \rr Pass@1 & \rr \textit{Favourable to the package.} Improvement in a package (S1-C1) and in ablations (S1133-C3); deteriorations with a step generated by a small model (S216-C2) and with knowledge generation (S2075-C2) & \rr approx. 20 (pkg. approx. 16; ingr. approx. 10) & \rr More calls; cost not measured; in a Java class, holistic generation outperformed the decomposed one (S365)\tabularnewline
\rr G8 & \rr Retrieved context (API, file, repository) & \rr Accessible repository or documentation & \rr Pass@1; accuracy; BLEU (S1951) & \rr \textit{Favourable; contrary in part.} Improvement in a package and in ablations (S357-C1; S852-C2); nulls and reversals (S806-C1; S1951-C1) & \rr approx. 45 (ingr. approx. 20) & \rr Oracle context is an upper bound (S670, S2145); more tokens; trained retrievers\tabularnewline
\rr G9 & \rr In-context examples & \rr Fixed, random or retrieved demonstrations & \rr Pass@1 & \rr \textit{Mixed.} Small differences of variable sign when examples are merely added (S432-C1; S591-C1); larger with trained selection (S170-C1); deterioration with faulty samples (S2075-C1) & \rr approx. 24 (ingr. approx. 14) & \rr Reversals by model and temperature; components outside the scope\tabularnewline
\rr G10 & \rr Specification and tests in the prompt & \rr Tests, pre- and postconditions or format rules in the prompt & \rr Accuracy; Pass@1; assertions passed & \rr \textit{Favourable in part.} Improvement with tests in the prompt (S1747-C1; S1086-C5) and with pre- and postconditions (S1812-C1); deterioration with Zero-Verbosity (S1606-C3) & \rr approx. 25 (ingr. approx. 9) & \rr Tests exposed in the prompt may also evaluate (S1613); S1812 not significant and aggregated\tabularnewline
\rr G11 & \rr Security instructions and constraints & \rr Prefixes, persona, CWE or guidelines & \rr Vulnerable samples by scanner; exposure to risk patterns; Pass@1 & \rr \textit{Mixed.} Deterioration and improvement, on distinct metrics (S604-C1; S397-C2; S1607-C1); lower Pass@1 with guidelines (S1195-C4). Do not recommend as an isolated measure & \rr 7 (pkg. and ingr.) & \rr Benchmark and scanner do not demonstrate security\tabularnewline
\rr G12 & \rr Wording, language, compression and automatic optimization & \rr Wording, natural language, pruning or search of the prompt & \rr Pass@1; incorrect solutions & \rr \textit{Mixed; insufficient per variant.} Rewordings worsen or do not change (S2162-C1; S323-C2); a language other than English worsens (S495-C1); guided pruning improves (S1644-C1). Do not recommend without a local test & \rr approx. 20 (ingr. approx. 8) & \rr One model or one study per variant; optimization in few studies (S650; S2053)\tabularnewline
\rr G13 & \rr Distribution of a fixed budget of calls (all) & \rr Choosing between repairing, generating new candidates or planning, with a limit of calls & \rr Pass 1@3; problems solved & \rr \textit{Insufficient.} The corpus does not guide how to split the calls. In the conditions extracted from the two studies with the budget of attempts equalized, the feedback loop did not outperform sampling (S479-C2; S645-C1); S479 reports gains with other models and budgets & \rr 2 studies & \rr 8B models and GPT-3.5 on competitive problems; they equalize attempts, not tokens or cost\tabularnewline
\rr G14 & \rr Domain languages and small local models & \rr E.g.: Solidity; models up to 8B & \rr Compilation; CodeBLEU (S556); Pass@1 & \rr \textit{Insufficient.} No contrast of a prompting technique in Solidity (S2055-C1 is cumulative; S556 only measures CodeBLEU). With small models, rows G3 and G6 record deteriorations (S759-C1; S1121-C1) & \rr no count of its own & \rr The evidence of the other rows comes from other models and languages\tabularnewline
\hline
\end{tabular}
\par\smallskip
\parbox{\textwidth}{\scriptsize Evidence: favourable, contrary or mixed, according to the direction of the comparisons; indirect, when it comes from another condition, domain or metric; insufficient, when the corpus has no contrast for the condition. Strength: comparisons and studies, without a score. Extraction by AI, checked by sample; syntheses of the 302 records used in v1.}
\end{table*}

\begin{table*}[!tp]
\caption{Decision-support table, version 2: debugging (D), translation (T) and refactoring (R)}\label{tab:apoiob}
\centering
\fontsize{6.6}{7.4}\selectfont
\renewcommand{\arraystretch}{0.95}
\setlength{\tabcolsep}{2.5pt}
\renewcommand*{\rr}{\raggedright\hspace{0pt}}
\begin{tabular}{p{0.55cm}p{2.0cm}p{2.4cm}p{1.9cm}p{5.6cm}p{1.6cm}p{3.1cm}}
\hline
\rr\textbf{ID} & \rr\textbf{Grouping} & \rr\textbf{Condition and verification} & \rr\textbf{Measured result} & \rr\textbf{Evidence and indication} & \rr\textbf{Strength} & \rr\textbf{Caveats}\tabularnewline
\hline
\rr D1 & \rr Test or execution feedback in a loop & \rr Tests that reproduce the defect, debugger or counterexample & \rr Correct fixes; plausible patches (S62) & \rr \textit{Favourable.} Improvement in most (S62-C2; S412-C3); tie with a 1B model (S114-C2) & \rr 13 (pkg. 8; ingr. 4) & \rr Assumes a test that triggers the defect; calls not equalized; hints from the answer key (S65); drop on Codeforces (S2164); a plausible patch is not a verified fix\tabularnewline
\rr D2 & \rr Diagnostic information in the prompt & \rr Failing test with error message, test output or defect report & \rr Correct fixes; repair rate; bugs localized (S1678) & \rr \textit{Favourable; contrary in part.} Improvement with triggering test and message (S647-C2), defect report (S647-C3) and test output (S1674-C5); deterioration with PDG (S1674-C3) & \rr 20 (ingr. in at least 8) & \rr Measured with perfect localization in S647; the deterioration with stack trace (S1678-C1) is in the localization of the defect, not in the repair; varies with the type of bug\tabularnewline
\rr D3 & \rr Retrieved context and ingredients & \rr Repository, history or base of patterns; in general with perfect localization or triggering tests & \rr Correct fixes; repair rate; exact match & \rr \textit{Favourable to the package.} Improvement in a package and in ablations (S275-C1; S435-C3; S2046-C2); injecting all the context at once worsens (S821-C2) & \rr 15 (ingr. 7) & \rr Perfect localization (S2046); triggering tests (S303); unequal budgets; different models between conditions\tabularnewline
\rr D4 & \rr Examples & \rr Fixed, historical or selected & \rr Exact match; programs fixed; overall score & \rr \textit{Mixed.} Improvement with selection or two examples (S25-C2; S53-C1); deterioration or tie with fixed or reference examples (S552-C3; S1598-C1). Do not recommend fixed examples & \rr 8 (ingr. in at least 4) & \rr Reversal between models (S2057); diverging tables (S25)\tabularnewline
\rr D5 & \rr Explicit reasoning & \rr CoT, ReAct or explanation of the error as the only mechanism added & \rr Repair rate; Pass@1; overall score & \rr \textit{Contrary.} Deterioration in 5 of 6 comparisons in which reasoning is the only mechanism added (S275-C2; S354-C2; S1429-C1; S1598-C2). Do not adopt as a default on the basis of this corpus; evaluate under the conditions of use & \rr 5 (ingr. 3; pkg. 2) & \rr The improvement (S412-C2) includes selection of examples; few contrasts, with a 1.3B model (S1598) and the quantum domain (S1429)\tabularnewline
\rr D6 & \rr Decomposition, planning and agents & \rr Several stages, tools or debate; in general with execution that reproduces the defect & \rr Correct fixes; repair rate & \rr \textit{Favourable to the package.} Improvement in a package (S810-C1; S699-C1) and in ablations (S810-C2; S1304-C3) & \rr 12 (pkg. 12; ablation in 4) & \rr Execution traces require reproduction (S1304-C2); different calls and models; reversals in subtotals (S303, S1304)\tabularnewline
\rr D7 & \rr Repair without reproduction and with unknown localization & \rr Only the defect report; the suite does not trigger the failure & \rr Correct fixes; plausible patches & \rr \textit{Insufficient.} There is no evidence to recommend the complete process. The defect report in the prompt is the only applicable favourable element, measured with perfect localization (S647-C3). Compilation and existing tests verify regressions, not the reported defect & \rr no count of its own & \rr Defects4J provides triggering tests; S303 and S2046 assume tests or localization\tabularnewline
\rr T1 & \rr Repair with tests and compiler & \rr Executable tests and compiler; iterative repair & \rr Computational accuracy; pass rate & \rr \textit{Favourable; contrary in one study.} Improvement in most: package with repair (S732-C1); tests in the prompt and repair (S329-C2; S329-C3). Drop in S2051-C1 (non-equivalent metrics); null at the ceiling (S919-C3) & \rr 11 (pkg. 9; ingr. 2) & \rr In S329, the tests in the prompt accompany most of the difference; S732 is a package with its own repair steps and limits; cost not measured\tabularnewline
\rr T2 & \rr Compiler-only feedback & \rr Compiler, without tests & \rr Compilation & \rr \textit{Insufficient.} Compiling does not demonstrate equivalence of behaviour; in S1163-C1, compilation stayed at the ceiling in both conditions & \rr 1 study & \rr No comparison measures equivalence with only the compiler\tabularnewline
\rr T3 & \rr Examples & \rr k examples; structural or retrieved selection & \rr Computational accuracy; BLEU; compilation & \rr \textit{Mixed.} Small improvement (S674-C1; S1032-C1) or deterioration (S964-C1). Do not recommend as a default & \rr 6 (ingr. 3; pkg. 3) & \rr BLEU does not measure correctness; the effect does not grow with k (S409)\tabularnewline
\rr T4 & \rr Explanation and two steps & \rr Explaining or representing before translating & \rr pass@1; rate of successful translation & \rr \textit{Contrary.} Deterioration (S193-C1; S836-C2; S836-C3); one-shot with reasoning improves (S836-C1). Do not adopt as a default on the basis of this corpus; evaluate under the conditions of use & \rr 2 (ingr. 1; pkg. 1) & \rr S836-C1 confounds example and reasoning; one model per study\tabularnewline
\rr T5 & \rr Legacy language or platform migration & \rr E.g.: COBOL; change of version or of platform & \rr — & \rr \textit{Insufficient.} No study of the corpus translates COBOL; migrations only appear without a counterfactual or in a cumulative contrast (S1220; S1803; S1836) & \rr no count of its own & \rr The absence of evidence does not demonstrate that the techniques are ineffective\tabularnewline
\rr R1 & \rr Refactoring type specified & \rr Expected type indicated in the prompt & \rr Recall in the identification of opportunities & \rr \textit{Indirect.} Recall from 15.6\% to 52.2\% (S511-C2) & \rr 1 study & \rr Measures the identification of opportunities, not the refactoring performed\tabularnewline
\rr R2 & \rr API knowledge: schema and changelog & \rr API documentation and changelog in the prompt & \rr Pass@1 & \rr \textit{Favourable in one study; indirect outside it.} Package from 4\% to 76\%; without the schema, 56\%; without the changelogs, 60\% (S2122-C1 to S2122-C3) & \rr 1 study & \rr Single study, on Terraform; number of calls not reported\tabularnewline
\rr R3 & \rr Examples & \rr Selected or one-shot examples; low baseline & \rr Passing the tests; Pass@1; correctness (S1822); complexity (S202) & \rr \textit{Favourable in pass rate; mixed in complexity.} Improvement (S202-C1; S4-C2; S1822-C1); S202 also measures complexity, with small differences that depend on the condition & \rr count of v1 for the three refactoring ingredients together: 11 (ingr. 6) & \rr Own sets; passing the tests does not imply reducing complexity; number of examples not established\tabularnewline
\rr R4 & \rr Feedback, iteration and reasoning & \rr Tests, compilation or detector; CoT & \rr Proportion that passes the tests; compilation; syntactic validity & \rr \textit{Insufficient.} Improvements (S1054-C1; S2183-C1); higher compilation without the feedback (S1685-C3); CoT with a small gain (S4-C1). Do not recommend as a default & \rr 7 (pkg. or ingr. in 6) & \rr Comparators with several ingredients; S739 without a baseline\tabularnewline
\rr R5 & \rr Repository scale and performance & \rr Interdependent changes in many files; performance requirement & \rr — & \rr \textit{Insufficient.} No study with a counterfactual at this scale; performance only appears against human code (S958) & \rr no count of its own & \rr Per-file checks do not establish the success of the complete migration\tabularnewline
\hline
\end{tabular}
\par\smallskip
\parbox{\textwidth}{\scriptsize Evidence: favourable, contrary or mixed, according to the direction of the comparisons; indirect, when it comes from another condition, domain or metric; insufficient, when the corpus has no contrast for the condition. Strength: comparisons and studies, without a score. Extraction by AI, checked by sample; syntheses of the 302 records used in v1.}
\end{table*}


\subsubsection{How to read the table and what it does not do}

The table is read by row: given the task and the condition of use, it summarizes the direction of the
extracted comparisons, the strength with which they were obtained and the caveats that prevent
generalization. ``Favourable'' indicates that most comparisons of the grouping recorded an
improvement on the metric of each study, and not that the technique will work in another project.
When the row is favourable to the package, the gain cannot be attributed to a single mechanism. The
table does not rank techniques: benchmarks, metrics and models differ between the rows, and
no average was computed. It also does not measure effectiveness in real use, since benchmarks do not
demonstrate productivity, security or cost reduction. The proportion of improvements among
the 805 comparisons (516 improvements, 145 deteriorations, 28 equal, 91 without direction and 25
descriptive) reflects the choice of the main comparison by the authors and publication
bias, and does not estimate an effect. Groupings with a single study, such as prompt
optimization and security requirements in translation (S650, S951), were left out of the table and
count as insufficient evidence. The two versions use the syntheses of 302 records; S1610, S1617, S667 and the studies of the last
batch of stage 3, extracted later, do not enter its rows. S433, cited in a row of v1, was later grouped
with S2053 for being a version of the same work (X1), and v2 cites S2053. Finally, the table reflects extraction by AI checked only by
sample: counts and directions may change with corrections outside it.
